% Generated draft: requires mathematical and source review.
\documentclass[11pt]{article}
\usepackage[margin=25mm]{geometry}
\usepackage{fontspec}
\setmainfont{Latin Modern Roman}
\usepackage{amsmath,amssymb,mathtools,unicode-math}
\setmathfont{Latin Modern Math}
\usepackage{enumitem,needspace,xurl,hyperref}
\setlength{\parindent}{0pt}
\setlength{\parskip}{5pt}
\setlength{\emergencystretch}{3em}
\newcommand{\N}{\mathbb{N}}
\newcommand{\Z}{\mathbb{Z}}
\newcommand{\Q}{\mathbb{Q}}
\newcommand{\R}{\mathbb{R}}
\newcommand{\C}{\mathbb{C}}
\newcommand{\F}{\mathbb{F}}
\newcommand{\A}{\mathbb{A}}
\newcommand{\PP}{\mathbb{P}}
\newcommand{\E}{\mathbb{E}}
\newcommand{\eps}{\varepsilon}
\newcommand{\dd}{\,\mathrm{d}}
\begin{document}
\section*{UnsolvedMath ID: 11000180}
\textbf{Source problem code:} catalogue entry
\subsection{Definitions and mathematical statement}
Let \(i:A\hookrightarrow B\) denote the inclusion in the source's Proposition 5. The principal problem is to prove that this inclusion is a bijection; since it is an inclusion, the remaining condition is
\[
 \forall b\in B\quad\exists a\in A:\quad i(a)=b.
\]
Neither Proposition 5 nor the definitions of \(A,B,i\) are supplied, so this logical form does not identify or replace the missing geometric assertion. The source expects additional analysis of the analytic and geometric foundations of the theory.

A further question concerns an asymptotic sequence \((f_k)\): given one element \(f_{k_0}\), can one describe \(f_l\) at other indices, beyond those covered by the Auroux--Katzarkov doubling formula? The extracted index is printed as ``\(2\ r k_0\)''; repeated doubling naturally suggests \(2^rk_0\), but the lost exponent must be checked. The objects \(f_k\), allowed indices and equivalence/description notion are absent. The final phrase ``Or perhaps better'' announces an unfinished reformulation and supplies no additional target.
\subsection{Short English statement}
Prove surjectivity of the missing Proposition 5 inclusion, and determine how one asymptotic-sequence element controls other indices beyond the stated doubling construction.
\subsection{Sources}
\begin{itemize}
\item \emph{Problems on Mapping Class Groups and Related Topics}, edited by Farb, Problem 2 and following incomplete question, supplied PDF page 177. \url{https://www.math.uchicago.edu/~farb/papers/mcgbook.pdf}. Recorded access 2026-07-29; extraction pdf-font-structure.
\item Proposition 5 and the Auroux--Katzarkov doubling formula are named, but their statements and a complete formula reference are absent.
\item Local UnsolvedMath registry, \texttt{data/unsolvedmath/problems.json}, permanent ID \(11000180\), exact source code \texttt{catalogue entry}; supplied OPEN-TRIAGE summary reports unverified status.
\end{itemize}
\end{document}
